\tikzset{
  edu usecase diagram/.style={
    font=\footnotesize,
    every node/.style={align=center}
  },
  edu system boundary/.style={
    draw,
    thick,
    rounded corners=4pt
  },
  edu usecase/.style={
    ellipse,
    draw,
    thick,
    fill=gray!8,
    minimum height=10mm,
    text width=29mm,
    inner sep=2pt
  },
  edu actor link/.style={
    draw,
    thick
  },
  edu decision wide/.style={
    edu decision,
    text width=34mm,
    aspect=1.8
  }
}

\newsavebox{\EduUseCaseBox}

\newcommand{\EduUseCaseFigure}[3]{%
  \begin{figure}[htbp]
    \centering
    \sbox{\EduUseCaseBox}{#1}%
    \ifdim\wd\EduUseCaseBox>0.96\textwidth
      \resizebox{0.96\textwidth}{!}{\usebox{\EduUseCaseBox}}%
    \else\ifdim\dimexpr\ht\EduUseCaseBox+\dp\EduUseCaseBox\relax>0.68\textheight
      \resizebox{!}{0.68\textheight}{\usebox{\EduUseCaseBox}}%
    \else
      \usebox{\EduUseCaseBox}%
    \fi\fi
    \caption{#2}
    \label{#3}
  \end{figure}
}

\newcommand{\EduActor}[2]{%
  \begin{scope}[shift={#1}]
    \draw[thick] (0,0.55) circle [radius=0.24];
    \draw[thick] (0,0.31) -- (0,-0.52);
    \draw[thick] (-0.42,0.02) -- (0,0.18) -- (0.42,0.02);
    \draw[thick] (0,-0.52) -- (-0.36,-1.02);
    \draw[thick] (0,-0.52) -- (0.36,-1.02);
    \node at (0,-1.48) {#2};
  \end{scope}
}
